% !PRINTED 443
\weberpage{443}
This work is considerably lightened by the fact that, from the consideration of the congruence group $L_7$, the degrees and cycles of the group are already known to us (\S.\,68). According to this there are $p^2-1=48$ elements of degree $7$, $A_7$, $\frac14 p(p-3)(p+1)=56$ elements of degree $3$, $A_3$, and $\frac14 p(p-1)^2=63$ elements of degree $4$ or $2$, $A_4, A_2$, in $G_{168}$.

The elements $A_7$ arrange themselves (with the inclusion of the identical substitution) into eight cycles of seven terms, the $A_3$ into 28 cycles of three elements, and the $A_4$ and $A_2$ into 21 cycles of four terms. Each of these last cycles contains one element $A_2$ and two elements $A_4$, and so, collecting these together, we obtain:
\[
\begin{aligned}
&48 \text{ elements of the } 7 \text{th degree}\\
&56 \text{ „} \quad 3 \text{rd „}\\
&42 \text{ „} \quad 4 \text{th „}\\
&21 \text{ „} \quad 2 \text{nd „}.
\end{aligned}
\]
If a point remains unchanged by $\nu$ substitutions of the group (the identical one included), then we shall call such a point a \gesperrt{$\nu$-fold pole} (\S.\,52).

For establishing the poles and axes it suffices to consider one cycle of each kind, since from it the others are derived by transformation (\S.\,69).

If we take the substitution $\tau$ of degree $7$ [\S.\,115, (4)], we find the vertices of the coordinate triangle as three separate poles. All powers of $\tau$ whose exponents are not divisible by 7 have the same three poles.

By applying the substitutions $\chi,\chi^2$ to the coordinates of these three poles, they pass cyclically into one another; but if one applies the substitutions $\Theta,\Theta^2,\Theta^3,\omega,\omega\Theta,\omega\Theta^2,\omega\Theta^3$, which are completely formed in \S.\,115, (5), (10), (16), (18), then the whole coordinate triangle passes into one entirely different from it (e.g.\ by $\omega$ into the triangle with the vertex coordinates $\alpha,\beta,\gamma$; $\beta,\gamma,\alpha$; $\gamma,\alpha,\beta$), and from this it is to be concluded that these poles are not more than seven-fold. There are eight triples of such seven-fold poles, all of which are different from one another, and each of them can be transformed into every other by a substitution of the group $G_{168}$.

We thus have the theorem:
%CONT>>
